\section{What telescopes: an amortized re-solving lemma}
\label{sec:amortized}

Although the number of expansions can be large, the continuous optimization
work needed to adapt to them admits a clean amortization.

For a safe point $0\leq x\leq x_\rho^\star$, define the one-sided violation
\begin{equation}
  \nu(x):=\norm{D^{-1/2}(\kappa_\rho(x))_-}_\infty.
  \label{eq:nu}
\end{equation}

\begin{lemma}[Violation to objective gap]
\label{lem:violation-gap}
For every safe point,
\begin{align}
 \norm{D^{-1/2}(x_\rho^\star-x)}_\infty
 &\leq\frac{\nu(x)}\alpha,
 \label{eq:nu-error}\\
 \nu(x)\leq\alpha\sigma
 \quad&\Longrightarrow\quad
 G_\rho(x)-G_\rho(x_\rho^\star)
 \leq\frac{\sigma^2}{2\rho}.
 \label{eq:nu-gap}
\end{align}
\end{lemma}

\begin{proof}
The proof of \cref{lem:one-sided-certificate} with
$\tau=\nu(x)/\alpha$ gives \cref{eq:nu-error}.  Put
$e=x_\rho^\star-x$.  Because $e$ is supported on $S_\rho$ and complementarity
annihilates the linear term,
\begin{equation}
 G_\rho(x)-G_\rho(x_\rho^\star)=\frac12e^\top Qe
 \leq\frac12\norm{e}_2^2.
\end{equation}
If $\nu(x)\leq\alpha\sigma$, then $e_i\leq\sigma\sqrt{d_i}$ by
\cref{eq:nu-error}.  Hence
\[
 \norm{e}_2^2
 \leq\sigma^2\vol(S_\rho)
 \leq\frac{\sigma^2}\rho,
\]
using \cref{eq:rppr-volume}.
\end{proof}

\begin{assumption}[Restricted accelerated contraction]
\label{ass:contraction}
On every fixed restricted QP, $k$ full restricted iterations reduce objective
error by a factor at most $\exp(-k/H)$, where
$H=C/\sqrt\alpha$ and $C$ is universal.
\end{assumption}

Strongly convex accelerated proximal gradient has this fixed-subspace
property.  The assumption says nothing about continuing momentum across a
subspace change.

\begin{theorem}[Amortized re-solving]
\label{thm:amortized}
Let $\sigma_{j-1}=2\sigma_j$ be geometric accuracy epochs and put
\begin{equation}
  E_j:=\frac{\sigma_j^2}{2\rho}.
\end{equation}
Suppose epoch $j$ begins at a safe point with
$\nu(x)\leq\alpha\sigma_{j-1}$, and restore restricted objective error to
$E_j$ after every support expansion.  Under \cref{ass:contraction}, the total
number of restricted iterations over all epochs down to $\sigma_J=\tau$ is
\begin{equation}
 K
 \leq\bigO\!\left(
       \frac1{\sqrt\alpha}\log\frac{\sigma_0}{\tau}
      \right)+N_{\mathrm{exp}},
 \label{eq:amortized-iterations}
\end{equation}
where $N_{\mathrm{exp}}$ is the number of nonempty expansion events.  Since
$\vol(U)\leq1/\rho$, full-sweep work obeys
\begin{equation}
 \mathcal W
 \leq\bigO\!\left(
       \frac1{\rho\sqrt\alpha}\log\frac{\sigma_0}{\tau}
      \right)
      +\bigO\!\left(\frac{N_{\mathrm{exp}}}{\rho}\right).
 \label{eq:amortized-work}
\end{equation}
\end{theorem}

\begin{proof}
By \cref{lem:violation-gap}, the global objective gap at the start of epoch
$j$ is at most $4E_j$.  Restricted optimal values decrease monotonically
toward the global value, so all expansion gains $\Delta_{j,\ell}$ in that
epoch telescope:
\begin{equation}
  \sum_\ell\Delta_{j,\ell}\leq4E_j.
  \label{eq:shock-telescope}
\end{equation}
If the pre-expansion restricted error is at most $E_j$, then immediately after
an optimum shift of size $\Delta_{j,\ell}$ it is at most
$E_j+\Delta_{j,\ell}$.  By \cref{ass:contraction}, restoring it to $E_j$
takes at most
\begin{equation}
 k_{j,\ell}
 =\left\lceil H\log\left(1+\frac{\Delta_{j,\ell}}{E_j}\right)\right\rceil.
\end{equation}
Therefore
\begin{align*}
 \sum_\ell k_{j,\ell}
 &\leq H\sum_\ell
       \log\left(1+\frac{\Delta_{j,\ell}}{E_j}\right)
       +N_{\mathrm{exp},j}\\
 &\leq\frac{H}{E_j}\sum_\ell\Delta_{j,\ell}
       +N_{\mathrm{exp},j}
 \leq4H+N_{\mathrm{exp},j}.
\end{align*}
Summing over $\bigO(\log(\sigma_0/\tau))$ epochs proves
\cref{eq:amortized-iterations}.  Multiplying full iterations by the uniform
volume cap proves \cref{eq:amortized-work}.
\end{proof}

This theorem improves over independently re-solving every enlarged subspace:
all continuous shock sizes telescope.  The unresolved term is discrete.  In
the worst case $N_{\mathrm{exp}}$ may be as large as the support size, and the
proof above charges a full current-support scan to every such event.
